\section{Necessary costs of the current graded standard-support bounds}
\label{sec:current-graded}

We now use the current ePrint's graded challenge count and squarefree
reconstruction formulas~\cite[Equations (54)--(58), (60), (63)]{DKTcurrent}.
For the standard derivative-capped, total-degree-capped source, their
quadratic list-budget and fourth-power regular-MCA-budget dependence
cannot be improved by choosing different cap parameters or a smaller
height certified by the same graded row test. This is a statement about
these numerical bounds, not actual component degrees, lists or exceptions.
The support restriction below is essential.

\begin{theorem}[Current graded standard-support costs]\label{thm:current-graded}
Let $N\ge12$ be divisible by four, $D=N/4-1$, and let $D+1<A\le N$
be an integer. For integers $m\ge1$ and $0\le b\le B$, use all monomials
\[
 X^xY_0^uY_1^v,\qquad
 v\le b,\quad u+v\le B,\quad x+Du+(D-1)v<mA,
 \qquad x,u,v\ge0.
\]
Assume characteristic zero or greater than every retained total jet degree.
Write $G_q,R_q$ for the exact source dimension and local rank on total jet
degree $q$, and put
\[
 a_* =\frac{AN-8}{N(N-8)},\qquad
 \varepsilon_*=a_*-a_0,\qquad c_0=\frac{1-2a_0}{8}.
\]
Suppose $0<\varepsilon_*\le1/2000$ and an integer $H\ge0$ passes the
graded row test for an affine received line:
\begin{equation}\label{eq:current-graded-test}
 \sum_q(H-q+1)_+(G_q-NR_q)>0.
\end{equation}
Then
\begin{equation}\label{eq:current-graded-costs}
 m>\frac{c_0}{\varepsilon_*},\qquad B>\frac m4,\qquad
 b\ge\frac m8,\qquad
 H+1>\frac{m}{36864\varepsilon_*},\qquad
 H>\frac{m}{73728\varepsilon_*}.
\end{equation}
The current squarefree list budget is therefore
$\Omega(D\varepsilon_*^{-2})$, and its regular-family MCA budget is
$\Omega(D^2\varepsilon_*^{-4})$, even after optimizing the latter's
integer incidence threshold.
\end{theorem}
\begin{proof}
Normalize the prefix surpluses by $N$:
\[
 P_j=\sum_{q\le j}(G_q/N-R_q).
\]
The left side of \eqref{eq:current-graded-test}, divided by $N$, is
$\sum_{j=0}^H P_j$. Some prefix is therefore positive. Total-degree
truncation preserves the full coefficient prefixes and downward-$Y_0$
closure. Applying the exact finite-length parameter-cost corollary of
Section~\ref{sec:gap-costs} to that prefix gives
$m>c_0/\varepsilon_*$, maximum $Y_0$ degree greater than $m/4$, and
maximum $Y_1$ degree greater than $m/4-1$. Hence $m\ge16$,
$B>m/4$, $b\ge m/8$, and $H>m/4-1\ge m/8$.

For every prefix, whether positive or negative, the reduction in
Section~\ref{sec:finite-length} deletes final diagonals of nonpositive
surplus and produces a saturated support $S_0$. Its exact coefficient
benefit is dominated by the leading benefit at $a_*$, with every retained
benefit positive. Thus each retained total degree is less than
$4ma_*<2m$, and $|S_0|\le2m^2+m\le3m^2$. The critical support converse
at $a_0$ yields
\begin{equation}\label{eq:graded-prefix-upper}
 P_j\le B_{m,1/4,a_*}(S_0)-r_m(S_0)
 \le m\varepsilon_*|S_0|\le3\varepsilon_*m^3.
\end{equation}
The middle inequality remains valid if some benefits at $a_0$ are
nonpositive, by deleting those monomials. Empty $S_0$ causes no problem.

Put $J=\lfloor m/8\rfloor$. The standard support contains all $q+1$
jet monomials on every degree $q\le J$. Their coefficient prefixes are
nonempty and locally saturated, since
\[
 L_q=mA-(D-1)q\ge m\bigl(A-(D-1)/8\bigr)\ge m.
\]
The exact rank formula, restricted to $\ell=q,\ldots,m-q-1$, gives
\[
 R_q=\sum_{\ell=0}^{m-1}\min\{\min(q,\ell)+1,m-\ell\}
 \ge(m-2q)(q+1)\ge\frac{3m(q+1)}4.
\]
The exact coefficient comparison gives $G_q/N\le ma_*(q+1)$.
As $a_*<47/100$, we obtain
\[
 G_q/N-R_q\le-\frac{m(q+1)}4,
 \qquad P_j\le-\frac{m(j+1)(j+2)}8\quad(0\le j\le J).
\]
Since $H\ge J$, the negative prefix area in the graded sum is at least
\[
 -\sum_{j=0}^J P_j\ge\frac{m(J+1)(J+2)(J+3)}{24}
 >\frac{m^4}{12288}.
\]
Every positive prefix is at most $3\varepsilon_*m^3$ by
\eqref{eq:graded-prefix-upper}. Positivity of the entire sum implies
$3\varepsilon_*m^3(H+1)>m^4/12288$, proving the $H+1$ bound.
Since $H\ge m/8\ge2$, also $H\ge(H+1)/2$, proving the last inequality
in \eqref{eq:current-graded-costs}.

For the current reconstruction formulas, write
\[
 \tau=2D-3,\quad u=1+\tau(B-1),\quad
 v=\min\{u,\tau(b-1)+D\},\quad F=Bv+b(u-v).
\]
We have $D\ge2$, $B\ge2$, $0\le v\le u$ and $\tau\ge D/2$. Thus
\[
 F=bu+(B-b)v\ge bu\ge\frac{DBb}{4}
 >\frac{Dm^2}{128}.
\]
The squarefree list budget is $\lambda_{\rm agr}F+S$, where
$\lambda_{\rm agr}=(N-D)/(A-D)\ge1$ and $S\ge0$. It is at least
$Dc_0^2/(128\varepsilon_*^2)$. The declared joint-degree bound is
\[
 J_{\rm reg}=H(2uv-v^2)+2(1+\tau H)F
 \ge\frac{D^2HBb}{4}
 >\frac{D^2c_0^3}{9437184\varepsilon_*^4}.
\]
For every permitted incidence threshold $D+1\le L\le A$, its multiplier
in the current regular MCA bound is
\[
 \frac{N-L+1}{A-L+1}\lambda_{\rm agr}\ge1.
\]
The other terms are nonnegative, completing the proof.
\end{proof}

\begin{remark}[Scope of current-method tightness]
The mandatory low-degree triangle is specific to the standard source
above; arbitrary downward supports may omit small derivative exponents.
No graded converse is proved here for those supports, for nonmonomial
sources, or for the current ePrint's alternative column-only test
(Equation (64)). The characteristic hypothesis remains the stronger
local-rank hypothesis of this note, not merely the reconstruction
condition $p>\max\{D,b\}$. Optimizing actual component degrees, exploiting
global kernel dependencies, or using additional equations is outside the
numerical ledger considered here. When $A/N=a_0+\varepsilon$ and
$\varepsilon\ge1/N$, the corrected margin $\varepsilon_*$ is
$\Theta(\varepsilon)$ as $N$ grows; near rounding one must retain its
explicit formula. No intrinsic lower bound on lists or MCA exceptions
follows from this theorem.
\end{remark}
